%%%%%%%% ICML 2026 EXAMPLE LATEX SUBMISSION FILE %%%%%%%%%%%%%%%%%

\documentclass{article}

% Recommended, but optional, packages for figures and better typesetting:
\usepackage{microtype}
\usepackage{graphicx}
\usepackage{subcaption}
\usepackage{booktabs} % for professional tables

% hyperref makes hyperlinks in the resulting PDF.
% If your build breaks (sometimes temporarily if a hyperlink spans a page)
% please comment out the following usepackage line and replace
% \usepackage{icml2026} with \usepackage[nohyperref]{icml2026} above.
\usepackage{hyperref}


% Attempt to make hyperref and algorithmic work together better:
\newcommand{\theHalgorithm}{\arabic{algorithm}}

% Use the following line for the initial blind version submitted for review:
\usepackage{icml2026}

% For preprint, use
% \usepackage[preprint]{icml2026}

% If accepted, instead use the following line for the camera-ready submission:
% \usepackage[accepted]{icml2026}

\usepackage{amsmath}
\usepackage{amssymb}
\usepackage{mathtools}
\usepackage{amsthm}

\usepackage{ytableau}
\usepackage{tikz}
\usetikzlibrary{arrows.meta}


% if you use cleveref..
\usepackage[capitalize,noabbrev]{cleveref}

%%%%%%%%%%%%%%%%%%%%%%%%%%%%%%%%
% THEOREMS
%%%%%%%%%%%%%%%%%%%%%%%%%%%%%%%%
\theoremstyle{plain}
\newtheorem{theorem}{Theorem}[section]
\newtheorem{proposition}[theorem]{Proposition}
\newtheorem{lemma}[theorem]{Lemma}
\newtheorem{corollary}[theorem]{Corollary}
\newtheorem{conjecture}[theorem]{Conjecture}
\theoremstyle{definition}
\newtheorem{definition}[theorem]{Definition}
\newtheorem{assumption}[theorem]{Assumption}
\newtheorem{observation}[theorem]{Observation}
\theoremstyle{remark}
\newtheorem{remark}[theorem]{Remark}


\newcommand\rows{\operatorname{rows}}
\newcommand\cols{\operatorname{cols}}


% Todonotes is useful during development; simply uncomment the next line
%    and comment out the line below the next line to turn off comments
%\usepackage[disable,textsize=tiny]{todonotes}
\usepackage[textsize=tiny]{todonotes}

% The \icmltitle you define below is probably too long as a header.
% Therefore, a short form for the running title is supplied here:
\icmltitlerunning{Differentiable Counting via Convolution}

\begin{document}

\twocolumn[
    \icmltitle{Differentiable Counting for Latent Rule Discovery via Convolution}

  % It is OKAY to include author information, even for blind submissions: the
  % style file will automatically remove it for you unless you've provided
  % the [accepted] option to the icml2026 package.

  % List of affiliations: The first argument should be a (short) identifier you
  % will use later to specify author affiliations Academic affiliations
  % should list Department, University, City, Region, Country Industry
  % affiliations should list Company, City, Region, Country

  % You can specify symbols, otherwise they are numbered in order. Ideally, you
  % should not use this facility. Affiliations will be numbered in order of
  % appearance and this is the preferred way.
  \icmlsetsymbol{equal}{*}

  \begin{icmlauthorlist}
    \icmlauthor{Byung-Hak Hwang}{KIAS}
    \icmlauthor{Chul-hee Lee}{CMC}
    \icmlauthor{Chanho Min}{Ajou}
    %\icmlauthor{}{sch}
    % \icmlauthor{Firstname8 Lastname8}{sch}
    % \icmlauthor{Firstname8 Lastname8}{yyy,comp}
    %\icmlauthor{}{sch}
    %\icmlauthor{}{sch}
  \end{icmlauthorlist}

    \icmlaffiliation{KIAS}{School of Mathematics, Korea Institute for Advanced Study}
    \icmlaffiliation{CMC}{CMC}
    \icmlaffiliation{Ajou}{Ajou University}

  \icmlcorrespondingauthor{Chanho Min}{first1.last1@xxx.edu}
  % \icmlcorrespondingauthor{Firstname2 Lastname2}{first2.last2@www.uk}

  % You may provide any keywords that you find helpful for describing your
  % paper; these are used to populate the "keywords" metadata in the PDF but
  % will not be shown in the document
  \icmlkeywords{Weakly Supervised Learning, Attention Mechanisms, AI for Math, Multiple Instance Learning, Rule Discovery from Counts}

  \vskip 0.3in
]

% this must go after the closing bracket ] following \twocolumn[ ...

% This command actually creates the footnote in the first column listing the
% affiliations and the copyright notice. The command takes one argument, which
% is text to display at the start of the footnote. The \icmlEqualContribution
% command is standard text for equal contribution. Remove it (just {}) if you
% do not need this facility.

% Use ONE of the following lines. DO NOT remove the command.
% If you have no special notice, KEEP empty braces:
\printAffiliationsAndNotice{}  % no special notice (required even if empty)
% Or, if applicable, use the standard equal contribution text:
% \printAffiliationsAndNotice{\icmlEqualContribution}

\begin{abstract}
    We study \textit{Latent Counting}, where only an aggregate count $Y=\sum_i z_i$ is observed for a bag of unlabeled instances, and the goal is to learn a latent instance-level selection rule that extrapolates to larger, unseen bags.
    Our approach combines a differentiable \textbf{Probabilistic Convolutional Aggregator} that computes the exact probability mass function of $Y$ with a \textbf{Purely-Relative Transformer} that removes absolute features and conditions attention only on learnable pairwise relations.
    We demonstrate our approach on a challenging rule-discovery problem in mathematics, the conjectural $(q,t)$-Littlewood--Richardson formula, achieving $>98\%$ accuracy while extrapolating beyond the training regime.
    More broadly, these results suggest that explicitly modeling integer-valued aggregation can support rule discovery under weak supervision.
\end{abstract}

	\section{Introduction}
	
	\paragraph{The Challenge of Blind Counting.} 
	Scientific discovery often necessitates inferring the discrete, atomic rules that govern complex systems based solely on aggregate observations. In fields ranging from particle physics to computational biology, researchers frequently observe a global property—such as a net charge or a total cell count—while the individual components contributing to that sum remain unlabeled. This setting gives rise to the problem of \textit{bag-level integer summation}: given only a total count, one must infer the latent selection rule determining which instances in the bag contribute to the sum. 
	
	\paragraph{Limitations of current MIL.} 
	This task is formally an instance of Multiple Instance Learning (MIL). However, standard MIL aggregation functions are ill-suited for exact integer inference. 
    \texttt{Max-pooling}, designed for classification, detects only the \textit{existence} of a signal (0 vs. 1) and fails to capture multiplicity.
    \texttt{Mean}/\texttt{sum}-pooling, typically used for regression, treat the target as continuous, blurring discrete structure.
    More fundamentally, such aggregation schemes do not model the probability mass function (PMF) of the aggregate itself, and therefore cannot represent exact integer-valued distributions.

\paragraph{Proposed Method: Convolutional MIL with Purely-Relative Representation.}
To reconcile differentiable learning with exact integer aggregation and robust extrapolation, we propose a two-part framework: \textbf{Probabilistic Convolutional Aggregator} and \textbf{Purely-Relative Transformer}.
The Probabilistic Convolutional Aggregator replaces mean/sum pooling by modeling each instance selection as a Bernoulli random variable and computing the probability distribution of the integer aggregates via convolution.
The Purely-Relative Transformer then promotes scale-robust generalization by removing absolute features and conditioning attention only on learnable pairwise relative relations between tokens, discouraging shortcut learning tied to absolute coordinates and encouraging structure-consistent rule extrapolation beyond the training regime.


    
	% \paragraph{Proposed Method: Convolutional MIL.} 
	% To reconcile differentiable learning with exact integer aggregation, we propose \textbf{Probabilistic Convolutional Aggregator}. Rather than summing expected values, we treat the latent selection of each instance as a Bernoulli random variable and compute the exact distribution of their sum. Since the PMF of a sum is the convolution of per-instance PMFs, we can efficiently compute the likelihood of an observed integer aggregate. This allows us to learn sharp, discrete selection rules within a fully differentiable framework, without resorting to continuous approximations.
	
	\paragraph{Case Study: The $(q,t)$-Littlewood--Richardson Problem.} 
	We validate our framework on an important open problem in mathematics: a conjectural factorization formula for the $(q,t)$-Littlewood--Richardson (LR) coefficients arising in the theory of Macdonald polynomials.
    We recast recovery of the signed exponents in the factorization as a latent signed counting problem.
    This perspective fits naturally within our probabilistic aggregation framework.
	
	\paragraph{Contributions.} 
    We contribute in three ways:
	\begin{itemize}
		\item \textbf{Methodological:} We introduce a generic \textit{Probabilistic Counting Wrapper} that enables standard neural backbones to model exact integer-supported PMFs via differentiable convolution.
		\item \textbf{Architectural:} We propose the \textit{Purely-Relative Transformer}, a backbone that enforces translational symmetry and locality, facilitating generalization beyond the training size regime.
		\item \textbf{Scientific:} We achieve high exact-match performance on the $(q,t)$-LR problem and demonstrate that the model recovers the classical Pieri rule as a \textit{sparse subset} of its learned selection patterns.
	\end{itemize}
	
	\section{Related Work}
	\label{sec:related_work}
	
	\subsection{Machine Learning for Mathematical Discovery}
	Recent work applies deep learning to pure mathematics, from symbolic manipulation to pattern discovery and hypothesis generation. Early work by \cite{Lample2020Deep} demonstrated that Transformers could solve integration and differential equations by treating them as translation tasks.
    More recently, AI has guided human intuition in open problems, uncovering novel relationships between knot invariants \cite{davies2021advancing} and discovering graph-theoretic counterexamples via reinforcement learning \cite{wagner2021constructions}.

    We study a weakly supervised variant to infer latent instance selections from aggregate integer supervision, requiring discrete, rule-like structure rather than smooth regression.

	
    \subsection{Multiple Instance Learning (MIL)}
    Our formulation falls within the Multiple Instance Learning (MIL) paradigm, where each training example is a bag of instances annotated only with a bag-level label.
    Early MIL formulations typically treated bag labels as binary
    indicators of the existence of at least one positive instance, a setting that
    motivated max/mean pooling and related aggregation rules~\cite{dietterich1997solving,maron1998multiple}.
    Modern set-function architectures (e.g. Deep Sets) justify sum/mean pooling as 
    universal approximators for permutation-invariant functions, while attention-based 
    pooling has become popular for its empirical performance and interpretability~\cite{zaheer2017deep,ilse2018attention}.
    
    Standard MIL pooling operations (max/mean/attention) are designed for 
    classification or real-valued regression and output continuous scores; they do not 
    explicitly model the discrete distribution of an integer sum over instances.
    Shukla et al.~\yrcite{shukla2023unified} address discrete counts directly by 
    proposing a differentiable method. Their algorithm relies on dynamic programming to obtain exact count probabilities, and is evaluated on synthetic benchmarks derived from MNIST. 
    While this formulation is powerful, the DP-style computation is not easily parallelizable, motivating alternative discrete aggregation methods that remain GPU-friendly and transferable to domain-specific tasks.
    % While this count-probability formulation is powerful, the dynamic-programming 
    % style computation used in the paper is not presented as a simple, fully
    % GPU-parallel operation. This motivates approaches that obtain
    % bag-level counts via GPU-friendly aggregation and are validated on realistic,
    % domain-specific tasks rather than only on synthetic image benchmarks.
    	
	% \subsection{Transformers and Relative Positional Encodings}
	% To generalize beyond training configurations, the model must be invariant to absolute position. Standard Transformers \cite{vaswani2017attention} rely on absolute positional encodings, which struggle to extrapolate beyond the training sequence length \cite{press2022train}. 
	
	% Relative Positional Encodings (RPE) \cite{shaw2018self, dai2019transformer} address this by encoding the distance between tokens rather than their coordinates. In the visual domain, \cite{parmar2018image} introduced local attention windows for 2D images to capture translation equivariance. We extend these ideas to the domain of Young diagrams. Unlike the dense, rectangular grids in \cite{parmar2018image}, partitions are not ... and rectangular. Our Windowed Geometric Attention enforces a hard locality constraint similar to the receptive fields in CNNs.	

\subsection{Transformers and Relative Positional Encodings}
Transformers, built on self-attention, are a general-purpose architecture for sequence and vision tasks~\cite{vaswani2017attention}.
Because Transformers do not encode order or coordinates intrinsically, they require positional information; absolute encodings can hurt generalization to unseen configurations (e.g., longer sequences), as highlighted in work on length extrapolation~\cite{press2022train}.

Relative positional encodings (RPE) instead model pairwise relations (e.g., distances) rather than absolute coordinates, improving performance and robustness in language and other sequence tasks~\cite{shaw2018self,dai2019transformer}.
In vision,
local attention / windowing schemes (e.g., Image Transformer, Swin Transformer)
combine spatial locality with attention on 2D grids~\cite{parmar2018image,liu2021swin}.
Because many RPE and windowing designs assume regularly indexed 1D sequences
or dense rectangular 2D patch grids, applying them to more flexible or irregular
coordinate sets requires additional design choices (for example,
how distances are measured or neighborhoods defined), which can limit
applicability to non-rectangular or structured domains.
	
\section{Problem Setup and Learning Objectives}
\label{sec:problem_formulation}

We formalize the task of inferring instance-level structure from aggregate integer-valued supervision as a weakly supervised learning problem. We consider two variations: latent counting and latent signed counting, where positive and negative contributions may cancel.

\subsection{Problem 1: Latent Counting as the Base Case}
Consider a dataset where each input is a bag $\mathcal{B} = \{x_1, \dots, x_N\}$ of $N$ unordered instances. Each instance $x_i$ has a latent binary selection variable $z_i \in \{0,1\}$ indicating whether it contributes to the observed aggregate, and $z_i$ is never observed.

The only supervision provided is the aggregate count $Y \in \mathbb{N}$, representing the total number of selected instances:
\[
    Y = \sum_{i=1}^N z_i.
\]
The goal is to learn $p_i=P(z_i=1)$ so that the induced sum distribution assigns high likelihood to the observed $Y$.

\subsection{Problem 2: Latent Signed Counting}
In many physical and mathematical contexts, contributions can be signed, such as charge, spin, or inclusion--exclusion from numerator--denominator structure, leading to cancellation. We extend the base formulation by associating each instance $x_i$ with a fixed known sign $s_i \in \{-1,+1\}$.

The observed target is now the \textit{net integer sum} $Y \in \mathbb{Z}$:
\[
    Y = \sum_{i=1}^N s_i \cdot z_i.
\]
This introduces a fundamental challenge: \textbf{cancellation}. For example, $Y=0$ may arise from selecting nothing or from selecting offsetting positive/negative instances. More generally, multiple distinct latent selections can yield the same observed sum, making the inverse problem more ambiguous than unsigned counting.

    
\subsection{The Dual Objective: Prediction and Discovery}
In standard supervised learning, predicting $Y$ accurately is often sufficient. In our setting, we also seek an interpretable instance-level explanation of the observed aggregate. Accordingly, our goal is twofold:

\begin{itemize}
    \item \textbf{Predictive Goal:} Minimize the negative log-likelihood of the observed aggregate $Y$ under the distribution induced by the predicted instance-level selection probabilities.
    \item \textbf{Discovery Goal:} Infer instance-level selection probabilities $\hat{p}_i$ that provide a plausible and interpretable explanation of the observed aggregate.
\end{itemize}

Crucially, our goal extends beyond minimizing prediction error on $Y$. We aim to learn a selection rule that produces consistent, interpretable instance-level attributions explaining the observed aggregate.
	
	
	\section{A Mathematical Problem as MIL}
	\label{sec:background}
	
	To ground our generic formulation, we focus on a specific object in mathematics: the $(q,t)$-Littlewood--Richardson coefficients. This provides a perfect testbed for Latent Signed Counting.
	
	\subsection{Partitions and Their Cells}
	A partition $\lambda = (\lambda_1, \lambda_2, \dots, \lambda_k)$ is a sequence of non-increasing positive integers, visualized as a grid of cells (Figure~\ref{fig:partition}). We identify a partition with its corresponding diagram. We denote $|\lambda| := \lambda_1+\dots+\lambda_k$.
    Also, denote by $\rows(\lambda)$ (resp., $\cols(\lambda)$) the number of rows (resp., columns)
    of $\lambda$, i.e., $\rows(\lambda) = k$ and $\cols(\lambda)=\lambda_1$.
    For any cell $u = (r, c)$ in row $r$ and column $c$, we define two intrinsic features:
	\begin{itemize}
		\item \textbf{Arm length} $a_\lambda(u)$: number of cells to the right of $u$.
        \item \textbf{Leg length} $\ell_\lambda(u)$: number of cells below $u$.
	\end{itemize}
    \begin{figure}
        \[
        \begin{ytableau}
           {~} & {u} & {a} & {a} \\
           {~} & {\ell} & {~} \\
           {~}
        \end{ytableau}
        \]
        \caption{A partition $(4,3,1)$}
        \label{fig:partition}
    \end{figure}
    % We say that a cell $u$ is of shape $(a_\lambda(u),\ell_\lambda(u))$.
    The following rational function serves as an atomic building block:
    for nonnegative integers $a$ and $\ell$,
    \[
        b(a,\ell) := \frac{1-q^{a}t^{\ell+1}}{1-q^{a+1}t^{\ell}},
    \]
    and for a partition $\lambda$ and a cell $u\in\lambda$,
    \( b_\lambda(u; q,t) := b(a_\lambda(u), \ell_\lambda(u)) \).
    For example, consider the cell $u=(1,2)$ in the partition shown in Figure~\ref{fig:partition}.
    Its arm length $a_\lambda(u)$ and leg length $\ell_\lambda(u)$ are $2$ and $1$, respectively.
    The cells contributing to $a_\lambda(u)$ and $\ell_\lambda(u)$ are indicated by $a$ and $\ell$, respectively.
    Although the row and column indices and the arm and leg lengths of a cell $u$ depend on the partition containing $u$, we write them as $r_u$, $c_u$, $a_u$, and $\ell_u$ when the partition is clear from the context.
    
    \subsection{$(q,t)$-Littlewood--Richardson coefficients}
    Littlewood--Richardson (LR) coefficients $c^\nu_{\lambda,\mu}$ and their $(q,t)$-generalizations
    $f^\nu_{\lambda,\mu}(q,t)$ are important objects that arise in many areas of mathematics; see e.g. \cite{LittlewoodRichardson1934,Macdonald1995,KnutsonTao1999,Vakil2006}.
    % While we do not provide the definitions of $c^\nu_{\lambda,\mu}$ and $f^\nu_{\lambda,\mu}(q,t)$,
    % their basic properties are as follows:
    % \begin{itemize}
    %     \item They are indexed by triples of partitions $(\lambda,\mu,\nu)$ with $\lambda,\mu\subseteq\nu$
    %           and $|\lambda|+|\mu|=|\nu|$.
    %     \item Each $c^\nu_{\lambda,\mu}$ is a nonnegative integer, while $f^\nu_{\lambda,\mu}(q,t)$
    %           is a rational function in $q$ and $t$.
    %     \item $f^\nu_{\lambda,\mu}(0,0)=c^\nu_{\lambda,\mu}$.
    % \end{itemize}
    Each $c^\nu_{\lambda,\mu}$ is a nonnegative integer, while $f^\nu_{\lambda,\mu}(q,t)$ is a rational function in $q$ and $t$ such that $f^\nu_{\lambda,\mu}(0,0)=c^\nu_{\lambda,\mu}$.
    Thus $f^\nu_{\lambda,\mu}(q,t)$ refines the classical coefficients $c^\nu_{\lambda,\mu}$.
    While $c^\nu_{\lambda,\mu}$ is well understood, the structure of $f^\nu_{\lambda,\mu}(q,t)$ remains much less so.
    % While many properties of $c^\nu_{\lambda,\mu}$ are well understood and explicit formulas exist, the structure of $f^\nu_{\lambda,\mu}(q,t)$ is substantially more intricate and remains poorly understood in general, due to their greater algebraic and combinatorial complexity.

    Stanley proposed a conjectural structural description of $f^\nu_{\lambda,\mu}(q,t)$
    for those triples with $c^\nu_{\lambda,\mu}=1$ (Conjecture~8.5 in \cite{Stanley1989}), which we reformulate as follows:
    \begin{conjecture} \label{conj:stanley_conjecture}
        Let $\lambda,\mu,\nu$ be partitions with $c^\nu_{\lambda,\mu}=1$. Then there exist subsets
        $X^\nu_{\lambda,\mu}\subseteq\lambda$, $Y^\nu_{\lambda,\mu}\subseteq\mu$, and
        $Z^\nu_{\lambda,\mu}\subseteq\nu$ such that
        \begin{equation} \label{eq:f_nlm-form}
            f^\nu_{\lambda,\mu}(q,t) =
            \dfrac{\prod_{u\in Z^\nu_{\lambda,\mu}} b_\nu(u;q,t)}
                  {\prod_{u\in X^\nu_{\lambda,\mu}} b_\lambda(u;q,t)
                   \times\prod_{u\in Y^\nu_{\lambda,\mu}} b_\mu(u;q,t)}.
        \end{equation}
    \end{conjecture}
    Informally, the conjecture posits a latent selection rule on cells of $\lambda,\mu,\nu$ such that $f^\nu_{\lambda,\mu}(q,t)$ factors into local cell contributions determined by the selected subsets.
    Note that such selections $X^\nu_{\lambda,\mu}$, $Y^\nu_{\lambda,\mu}$ and $Z^\nu_{\lambda,\mu}$ may not be unique.

    The conjecture is still widely open, but some very special cases are known.
    The most notable case is the case where $\mu$ is a single row partition ($\rows(\mu)=1$).
    In this case, Macdonald gave an explicit description of the sets $X^\nu_{\lambda,\mu}$,
    $Y^\nu_{\lambda,\mu}$, and $Z^\nu_{\lambda,\mu}$; see p. 340 in \cite{Macdonald1995} for details.
    We illustrate an example of the description in Figure~\ref{fig:Pieri}.
    The rule for this specific case is called the \emph{Pieri rule}, which will play a role
    in training and testing our model.


    
    \begin{figure}
        \ytableausetup{centertableaux, smalltableaux}
        \[
        \lambda = \begin{ytableau}
           {~} & *(yellow!50){~} & *(yellow!50){~} & {~} \\
           {~} & *(yellow!50){~} & *(yellow!50){~} \\
           {~}
        \end{ytableau} \qquad
        \mu = \begin{ytableau}
           *(green!50){~} & *(green!50){~} & *(green!50){~}
        \end{ytableau} \qquad
        \nu = \begin{ytableau}
           {~} & *(blue!50){~} & *(blue!50){~} & {~} & *(blue!50){~} \\
           {~} & *(blue!50){~} & *(blue!50){~} \\
           {~} & *(blue!50){~} & *(blue!50){~}
        \end{ytableau}
        \]
        \caption{Example of the Pieri rule.
        Let $\lambda=(4,3,1)$, $\mu=(3)$, and $\nu=(5,3,3)$, then by the Pieri rule,
        the sets $X^\nu_{\lambda,\mu}$ (in yellow), $Y^\nu_{\lambda,\mu}$ (in green),
        and $Z^\nu_{\lambda,\mu}$ (in blue) give $f^\nu_{\lambda,\mu}$ in the sense of
        Eq.~\eqref{eq:f_nlm-form}.}
        \label{fig:Pieri}
    \end{figure}

    
    \subsection{Mapping to MIL}
    Since each atomic rational function $b_\lambda(u;q,t)$ is determined by only
    the arm and leg lengths $(a_u, \ell_u)$, some cancellations may appear
    in Eq.~\eqref{eq:f_nlm-form}. Hence we can rewrite $f^\nu_{\lambda,\mu}(q,t)$ as follows:
	\begin{observation}
		If $c^{\nu}_{\lambda,\mu}=1$, then $f^\nu_{\lambda,\mu}(q,t)$ factors
        into terms grouped by $(a,\ell)$:
		\[ \label{eq:factored_f}
			f^{\nu}_{\lambda,\mu}(q,t) = \prod_{(a,\ell)\in \mathbb{Z}_{\ge 0}^2} b(a,\ell)^{\beta_{\lambda, \mu, \nu, a,\ell}}.
		\]
		The exponent $\beta_{\lambda, \mu, \nu, a,\ell}$ represents the net signed count of selected cells sharing the value $(a,\ell)$.
        % More precisely, a cell in $\nu$ (resp., $\lambda$ or $\mu$) has the value $(a,\ell)$, then it contributes $+1$ (resp., $-1$) to $\beta_{\lambda, \mu, \nu, a,\ell}$.
        More precisely, each selected cell $u$ with $(a_u,\ell_u)=(a,\ell)$ contributes $+1$ if $u\in\nu$ and $-1$ if $u\in\lambda\sqcup\mu$.
        See Figure~\ref{fig:cell_selection} for the visualization of the observation.
	\end{observation}
	
	Each triple $(\lambda,\mu,\nu)$ yields finitely many bags indexed by arm--leg pairs $(a,\ell)$, and we summarize the associated notation below.
	\begin{itemize}
    
            \item \textbf{Index set $\mathcal S_{\lambda,\mu,\nu}$:}
                the set of distinct arm--leg pairs $(a_u,\ell_u)$ attained by cells $u\in\lambda\sqcup\mu\sqcup\nu$.
                % Unique collection of all $(a(u),\ell(u))$ for all $u\in \lambda\sqcup\mu\sqcup\nu$.
            \item \textbf{Bag $\mathcal{B}_{\lambda, \mu, \nu, a,\ell}$:}
                all cells $u\in \lambda\sqcup\mu\sqcup\nu$ having the same value $(a,\ell)$.
		\item \textbf{Label $\beta_{\lambda,\mu,\nu,a,\ell}$:}
            the observed exponent.
		\item \textbf{Latent Variable $z_u$:}
            an indicator of whether cell $u$ is selected, $z_u\in\{0,1\}$.
		\item \textbf{Sign $s_u$:}
            $+1$ if $u \in \nu$, $-1$ if $u \in \lambda \sqcup \mu$.
	\end{itemize}
	
    Our goal is to train a model such that, for every bag $(\lambda, \mu, \nu, a,\ell)$,
    \[
\beta_{\lambda,\mu,\nu,a,\ell}=\sum_{u\in\mathcal{B}_{\lambda, \mu, \nu, a,\ell}} s_u \cdot z_u .
    \]
    \input{./figure-selection_example.tex}
    % Figure~\ref{fig:cell_selection} depicts the 

\section{Methodology}
\paragraph{Overview.}
Our system consists of two modules: a selector backbone $f_\theta$ and Probabilistic Convolutional Aggregator $\mathcal{A}$.
Given a triplet $(\lambda,\mu,\nu)$, we treat all cells as tokens and apply the Purely-Relative Transformer in parallel to produce a selection probability $p_i$ for each token $u_i$.
From each $p_i$, we construct an atomic PMF $\mathbf d_i$, whose support depends on whether $u_i\in \nu$ or not.
The aggregator $\mathcal{A}$ maps the collection of atomic PMFs $\{\mathbf d_i\}$ to a family of exact bag-level distributions indexed by $(a,\ell)$:
\[
\mathcal{A}\!\left(\{\mathbf d_i\}\right)
=\Big(\hat{P}_{\lambda,\mu,\nu,a,\ell}\Big)_{(a,\ell)\in \mathcal S_{\lambda,\mu,\nu}},
\]
where $\hat{P}_{\lambda,\mu,\nu,a,\ell}$ denotes the model-estimated PMF of the bag exponent $\beta_{\lambda,\mu,\nu,a,\ell}$.
A point prediction is obtained by
\[
\widehat{\beta}_{\lambda,\mu,\nu,a,\ell}
:= \arg\max_{k\in\mathbb{Z}} \hat{P}_{\lambda,\mu,\nu,a,\ell}(k).
\]
This construction enables end-to-end maximum-likelihood training under aggregate-only supervision.
While $\mathcal{A}$ can wrap any backbone $f_\theta$, the purely-relative design of $f_\theta$ is crucial for extrapolation to larger grids.
Figure~\ref{fig:pipeline} illustrates the overview of the pipeline.

    \input{./figure-pipeline.tex}



	\subsection{Methodology I: Probabilistic Convolutional Aggregator}
	\label{sec:method_conv}

    We propose a differentiable aggregation layer that solves the \textit{Latent Counting} problem (Problem~1) and naturally extends to the \textit{Signed} case (Problem~2). 
    The layer can wrap an arbitrary backbone, mapping its instance-level selection probabilities to bag-level PMFs via convolution.

\subsubsection{The Latent Counting Framework}
\label{sec:universal_framework}

First, we consider the general case of summing independent binary random variables. Standard approaches often approximate the sum of the binary selections $Y=\sum_i z_i$ using expectations, i.e., $\hat{Y}=\sum_i p_i$. This reduces the entire uncertainty of the discrete random sum to a single scalar. For instance, if $N=10$ and $p_i=0.5$ for all $i$, the expectation surrogate yields $\hat{Y}=5$, whereas the true object of interest is the full PMF of $Y$, supported on $\{0,1,\dots,10\}$, i.e., an $11$-dimensional distribution.

Accordingly, we compute the full PMF of the sum. 
Let $p_i\in[0,1]$ denote the predicted selection probability for instance $i$.
The Bernoulli PMF of $z_i$ is represented as a length-2 vector
\[
\mathbf d_i = (1-p_i,\; p_i).
\]
Since the $z_i$ are independent, the PMF of their sum is obtained by discrete convolution of the individual PMFs:
\[
P \;=\; \mathbf d_1 * \mathbf d_2 * \cdots * \mathbf d_N.
\]


\paragraph{Relation to prior work.}
Shukla et al.~\cite{shukla2023unified} compute the exact count probability via a dynamic-programming recursion with time $O(ks)$ for a fixed sum and $O(k^2)$ to obtain all sums. 
Our aggregator produces the same PMF, but we adopt a convolutional viewpoint that maps the recursion to repeated 1D convolutions. 
This form is implementation-friendly in our setting, since each input induces many bags and the aggregations can be batched via grouped Conv1D kernels on GPUs.



	\subsubsection{Extension to Signed Domains}
	We generalize to \textit{Latent Signed Counting}, where each instance has a known sign $s_i\in\{-1,+1\}$ and the observed label is
	$Y=\sum_i s_i z_i\in\mathbb{Z}$.
	Each instance contributes either $0$ or $s_i$, hence the atomic PMF:
	\[
		\mathbf d_i=
		\begin{cases}
			(0,1-p_i, p_i) & s_i=1,\\
			(p_i, 1-p_i , 0)   & s_i=-1.
		\end{cases}
	\]
	For a bag with $K$ instances, the signed sum lies in $[-K,K]$.
	When $s_i=-1$, the mass at $k=-1$ corresponds to a left shift, allowing exact cancellation via the same convolution operator.
	Training minimizes the signed negative log-likelihood.
	
\subsubsection{Scalable implementation via grouped 1D convolution}
\label{sec:implementation}

The convolutional view yields a direct and efficient implementation.
For each bag, we start from a point-mass PMF at $0$ and iteratively convolve it with the atomic PMFs.
% \[
% P \;=\; \mathbf d_1 * \mathbf d_2 * \cdots * \mathbf d_K.
% \]
% where each $\mathbf{d_i}$ is a length-3 kernel.
Since each triplet $(\lambda,\mu,\nu)$ induces many bags indexed by $(a,\ell)$, we batch all bags and compute these convolutions in parallel using standard grouped 1D convolutions. 
This avoids per-instance Python loops and maps the entire aggregation step to a small number of GPU kernels.

Optionally, the sequential chain can be further accelerated by a balanced divide-and-conquer schedule:
\[
\big(\mathbf d_1 * \mathbf d_2 * \mathbf d_3 * \mathbf d_4\big)
\;=\;
\big(\mathbf d_1 * \mathbf d_2\big) * \big(\mathbf d_3 * \mathbf d_4\big),
\]
which reduces the convolution depth at the cost of intermediate kernels with larger support.
In our current regime, bags are small enough that the straightforward grouped-convolution implementation is sufficient, but this suggests a clear path to scaling (e.g., via balanced trees or FFT-based methods for larger bags).

	

	
\subsection{Methodology II: Purely-Relative Transformer}
\label{sec:method_selector}
While the Convolutional Aggregator (Section \ref{sec:method_conv}) handles exact integer arithmetic, the system's ability to discover valid laws depends on a ``Selector'' backbone $f_\theta$ identifying consistent substructures.
We address the \textit{Extrapolation Gap}:
while standard models achieve $>99.99\%$ accuracy on random splits, they fail on larger grids, indicating overfitting to absolute coordinates or features.
To bridge this gap, we introduce the \textbf{Purely-Relative Transformer}, a backbone specifically designed to rely exclusively on relative spatial and shape dependencies, ensuring invariance to grid size and absolute position.

\subsubsection{Input Representation}

\paragraph{Tokens.}
The input is the union of cells from three partitions $\lambda, \mu,\nu$.
Each cell $u\in \lambda\sqcup \mu\sqcup \nu$ is treated as a token. % $u \in \mathcal{V}$.
We associate each token with a source identifier $\mathrm{id}_u \in \{1,2,3\}$ indicating whether $u$ comes from $\lambda$, $\mu$, or $\nu$.
We also compute geometric quantities (row/column $(r_u,c_u)$ and arm/leg $(a_u,\ell_u)$) as metadata.

\paragraph{Standard Transformer (absolute).}
As a conventional baseline, we embed the raw scalars and add absolute positional encodings:
\begin{align*}
\mathbf{x}_u^{\mathrm{std}} &= [r_u,\ c_u,\ a_u,\ \ell_u,\ \mathrm{id}_u] \in \mathbb{R}^5, \\
h_{u,\mathrm{std}}^{(0)} &= \mathrm{Linear}(\mathbf{x}_u^{\mathrm{std}}) + \mathrm{PE}(r_u,c_u).
\end{align*}

\paragraph{Ours (ID-only input).}
To eliminate absolute geometric channels from the token features, our model uses an \emph{ID-only} input:
\begin{align*}
\mathbf{x}_u^{\mathrm{ours}} = [\mathrm{id}_u] \in \mathbb{R}^1,\quad 
h_{u,\mathrm{ours}}^{(0)} = \mathrm{EmbedID}(\mathrm{id}_u),
\end{align*}
where $\mathrm{EmbedID}(\cdot)$ is a learnable embedding (or equivalently a linear map applied to a one-hot encoding of $\mathrm{id}_u$).
Crucially, our model does \emph{not} provide $(r_u,c_u,a_u,\ell_u)$ to the token features and does not use absolute positional encodings.

\subsubsection{Relative Attention}

\paragraph{Standard attention.}
Standard Transformers inject absolute location information into token representations via positional encodings.
The attention logits are
\[
e^{\mathrm{std}}_{ij} \;=\; \frac{\langle Q_i, K_j\rangle}{\sqrt{d}}.
\]
% with $Q_i = h_i W_Q$ and $K_j = h_j W_K$.

\paragraph{Our model's attention.}
Our model removes standard positional encodings and uses ID-only token embeddings at initialization.
Geometric structure is injected \emph{only} as an additive relative bias on the attention logits:
\[
e^{\mathrm{ours}}_{ij} \;=\; \frac{\langle Q_i, K_j\rangle}{\sqrt{d}} \;+\; b_{\mathrm{rel}}(i,j).
\]

\paragraph{Relative bias.}
We parameterize $b_{\mathrm{rel}}(i,j)$ as a sum of learnable lookup tables over relative differences:
\[
b_{\mathrm{rel}}(i,j)=
b_{r}(\Delta r_{ij}) + b_{c}(\Delta c_{ij})+\; b_{a}(\Delta a_{ij}) + b_{\ell}(\Delta \ell_{ij}),
\]
where $\Delta r_{ij}=r_{u_i}-r_{u_j}$, $\Delta c_{ij}=c_{u_i}-c_{u_j}$,
$\Delta a_{ij}=a_{u_i}-a_{u_j}$, and $\Delta \ell_{ij}=\ell_{u_i}-\ell_{u_j}$.
Later, we optionally clamp these differences to a fixed window to support size extrapolation.
\subsection{Window Clamping}
\label{sec:clamping}

A key obstacle to extrapolation on larger grids is that relative offsets grow in magnitude: as the grid expands, $\Delta r$ and $\Delta c$ can take values never observed during training. Yet many mathematical rules depend primarily on \emph{local} distances and beyond a neighborhood, only coarse directional information is needed.

We encode this bias by applying window-clamping to all relative-difference inputs to a fixed range $[-K,K]$:
\[
\mathrm{clamp}(x;K) = \max(-K,\min(x,K)).
\]
Clamping maps every offset to $[-K,K]$, eliminating unseen magnitudes at inference time while preserving the sign ($-K$ vs.\ $+K$).

\section{Experiments}
\label{sec:experiments}
	
We evaluate whether, under aggregate-only supervision, the proposed Probabilistic Convolutional Aggregator and Purely-Relative Transformer enable a model to (i) extrapolate a consistent rule to out-of-distribution regimes beyond the training domain and (ii) infer coherent instance-level selections in identifiable cases.

% We emphasize out-of-distribution (OOD) generalization to larger, unseen configurations, and quantify rule coherence by comparing the predicted selected-cell set to known mathematical rules in Pieri cases.

\subsection{Dataset construction and splits}
\label{sec:dataset_setup}

\textbf{Training set (general + Pieri augmentation), $D_{\text{train}}$.}
Our core training distribution consists of $(\lambda,\mu,\nu)$ triplets generated by an external procedure, restricted to small shapes where computation is feasible:
\[
    \rows(\lambda),\rows(\mu) \le 4,\quad \cols(\lambda),\cols(\mu) \le 8,
\]
together with their corresponding $\nu$ satisfying $c^\nu_{\lambda,\mu}=1$.
However, the ``$\rows \le 4$'' regime is extremely limited and provides only a few examples of strip-like configurations that dominate known closed-form Pieri cases.
We therefore augment training with Pieri instances of size up to 6 (that is, $\rows(\lambda),\cols(\lambda)\le 6$ and $\mu=(k)$ where $k\le 6$).
Thus, $D_{\text{train}}$ combines (a) small-regime general instances ($\rows \le 4$, $\cols \le 8$) and (b) Pieri instances up to size 6.

\textbf{Test set 1 (general OOD test set), $D_{\text{test1}}$.}
To assess extrapolation beyond the training size regime, we construct a general evaluation set at larger configurations, denoted $D_{\text{test1}}$.
We generate general triplets with up to 6 rows and then remove any overlap with the training set to prevent leakage:
\( D_{\text{test1}} \cap D_{\text{train}} = \emptyset. \)

\textbf{Test set 2 (extreme OOD test set), $D_{\text{test2}}$.}
Ideally, one would also evaluate on \emph{general} triplets at even larger sizes (e.g., rows 8). 
In practice, generating general $(q,t)$-LR exponents at this scale is computationally prohibitive.
To still probe extreme extrapolation, we therefore construct a second OOD set, denoted $D_{\text{test2}}$, using Pieri cases at size up to 8 rows.
We again enforce de-duplication by excluding any Pieri 8 instances that appear in training.

\subsection{Evaluation Metrics}
\label{sec:metrics}

\textbf{(1) Bag-level aggregate accuracy.}
Each triplet $(\lambda,\mu,\nu)$ decomposes into multiple disjoint bags indexed by $ \mathcal{B}_{\lambda, \mu, \nu, a,\ell}$, each with an integer target $\beta_{\lambda,\mu,\nu,a,\ell}$.
We report the \emph{bag-level accuracy} as the fraction of bags whose predicted exponent matches the ground truth exactly:
\begin{align*}
&\mathrm{Acc}_{\mathrm{bag}}\\
&~~:= \frac{1}{M}\sum_{(\lambda,\mu,\nu)\in D}\ \sum_{(a,\ell)\in \mathcal S_{\lambda,\mu,\nu}}
\mathbb{I}\!\left(\widehat{\beta}_{\lambda,\mu,\nu,a,\ell}=\beta_{\lambda,\mu,\nu,a,\ell}\right),
\end{align*}
where $M$ is the total number of bags in dataset $D$:
\( M:= \sum_{(\lambda,\mu,\nu)\in D}\left|\mathcal S_{\lambda,\mu,\nu}\right|. \)


\textbf{(2) Instance-level rule coherence on Pieri cases.}
For Pieri instances, a known rule specifies a reference selected-cell set $\mathcal{Z}_{\text{Pieri}}$.
Our selector produces an instance-level probability $p_u \in [0,1]$ for each cell $u$, interpreted as the probability that $u$ is selected.
We convert these probabilities into a predicted selected-cell set using a fixed threshold,
\[
\hat{\mathcal{Z}} \;=\; \{u \in\lambda\sqcup\mu\sqcup\nu : p_u \ge 0.5\},
\]
and compare it to $\mathcal{Z}_{\text{Pieri}}$.
To obtain near-binary selection probabilities $p_u$, we add a small entropy regularizer during training.
As a result, thresholding at $0.5$ yields stable predicted sets with few near-threshold probabilities.
We quantify coherence using two indicators:
\[
\text{Inclusion} = \mathbb{I}\big(\hat{\mathcal{Z}} \subseteq \mathcal{Z}_{\text{Pieri}}\big), \quad
\text{ExactSet} = \mathbb{I}\big(\hat{\mathcal{Z}} = \mathcal{Z}_{\text{Pieri}}\big).
\]
We report the inclusion rate and exact-set-match rate over Pieri test samples.

\paragraph{Training details.}
Unless otherwise stated, we use a Transformer selector with 4 layers, 4 heads, and model width 128.
We train for 200 epochs with batch size 512 using Adam (learning rate $10^{-3}$).
During training, selection probabilities are computed as $p=\sigma(l/\tau)$ and we anneal the temperature from $\tau=10$ to $\tau=1$; we also add an entropy regularizer with coefficient $0.01$.
For relative-bias attention, we use window clamping with $K=2$, so the model distinguishes only coarse relative relations (e.g., adjacent vs.\ far, and direction) rather than exact long-range distances.
All test results are averaged over 5 random seeds; we report mean and standard deviation.




\section{Results}
\label{sec:results}

\subsection{Aggregate Counting Performance}
\label{sec:quant_results}

\begin{table}[h]
	\centering
\caption{\textbf{Counting performance (bag-level accuracy).}
A sum-pooling regression baseline trained with MSE performs poorly, whereas the proposed Probabilistic Convolutional Aggregator achieves near-perfect accuracy.}

	\label{tab:ablation_counting}
	\begin{tabular}{lcc}
		\toprule
		\textbf{Method} & 		$\mathrm{Acc}_{\text{bag}}(D_{\text{train}})$ &
		$\mathrm{Acc}_{\text{bag}}(D_{\text{test1}})$  \\
		\midrule
		Regression (MSE Loss) & 58.34 & 64.30 \\
		\textbf{Convolution (Ours)} & \textbf{99.99} & \textbf{98.47} \\
		\bottomrule
	\end{tabular}
\end{table}
\paragraph{Effect of PMF aggregation.}
We test whether aggregation benefits from explicitly modeling the full distribution of an integer-valued sum.
To control for the selector, both methods use the same backbone $f_\theta$ (the Purely-Relative Transformer described in Section~\ref{sec:method_selector}) and differ only in the aggregation objective.
As a continuous surrogate, we train a sum-pooling regressor that predicts $\hat{\beta}_{sum}=\sum_i p_i$ and minimizes an MSE loss.
In contrast, our Probabilistic Convolutional Aggregator constructs an exact PMF of the aggregate $\hat{P}$ and maximizes the log-likelihood of the observed integer target.
Table~\ref{tab:ablation_counting} shows that the MSE regressor achieves only $58.34\%$ bag-level accuracy on $D_{\text{train}}$, whereas the PMF-based objective reaches $99.99\%$ on $D_{\text{train}}$ and $98.47\%$ on $D_{\text{test1}}$.
This gap is consistent with the fact that our supervision consists of multiple signed bag-level counts, where a continuous regression surrogate discards distributional information needed for stable learning. 
Unlike the standard Transformer baseline, our Purely-Relative Transformer uses relative biases for both spatial offsets $(\Delta r,\Delta c)$ and intrinsic shape offsets $(\Delta a,\Delta \ell)$, which improves OOD extrapolation.






\subsection{Out-of-Distribution Generalization}
\label{sec:ood_results}
\begin{table}[t]
	\centering
\caption{\textbf{Out-of-distribution performance.}
Bag-level accuracy on two OOD splits, $D_{\text{test1}}$ and $D_{\text{test2}}$.}
	\label{tab:headline_main}
	\begin{tabular}{lcc}
		\toprule
		\textbf{Model} &
		$\mathrm{Acc}_{\text{bag}}(D_{\text{test1}})$ &
		$\mathrm{Acc}_{\text{bag}}(D_{\text{test2}})$ \\
		\midrule
		Transformer & 91.67(1.13) & 86.01(3.07) \\
		\textbf{Ours} & \textbf{98.47(0.35)} & \textbf{98.22(0.48)} \\
		\bottomrule
	\end{tabular}
\end{table}
In this task, in-distribution (ID) evaluation can be misleading: with the Convolutional Aggregator fixed, all compared backbones achieve saturated bag-level accuracy on an ID $8{:}2$ split of the training distribution ($\mathrm{Acc}_{\mathrm{bag}}\ge 99.99\%$).
The central challenge is therefore out-of-distribution (OOD) generalization to unseen and structurally different configurations, as is common in rule-learning settings.
We evaluate extrapolation on two OOD splits: $D_{\text{test1}}$ (general instances with $\rows \le 6$) and $D_{\text{test2}}$ (Pieri instances with $\rows,\cols \le 8$).
All models share the same Probabilistic Convolutional Aggregator, and we vary only the selector backbone to isolate the effect of architectural inductive biases on OOD transfer.

\paragraph{Out-of-distribution performance.}
Table~\ref{tab:headline_main} summarizes bag-level accuracy on two OOD splits.
Relative to a standard Transformer, our purely-relative backbone improves accuracy from $91.67\%\to 98.47\%$ on $D_{\text{test1}}$ and from $86.01\%\to 98.22\%$ on the harder extrapolation set $D_{\text{test2}}$.
Notably, our method also exhibits substantially lower run-to-run variability (standard deviations $0.35$ and $0.48$), suggesting that removing absolute positional shortcuts yields more stable optimization and transfer across OOD regimes.


\paragraph{Ablation: positional information as a shortcut.}
Table~\ref{tab:ablation_pe_hurt_ours} shows that any access to absolute position can hurt OOD transfer.
Even with relative attention biases, positional encodings introduce an implicit absolute-position channel that the model can exploit in the small training regime but cannot extrapolate reliably to larger grids.
Removing PE forces the backbone to rely on purely relative cues that remain well-defined across OOD regimes.



\paragraph{Ablation: relative spatial bias alone.}
We next ask whether adding relative spatial bias to attention (while keeping standard positional encodings) improves extrapolation.
Table~\ref{tab:ablation_rel_rc} shows that, under the standard Transformer backbone, switching from absolute attention to a relative $(r,c)$ bias does not consistently improve OOD accuracy.
This suggests that relative bias by itself is not the main driver of transfer when absolute positional signals are still available.


\paragraph{Summary.}
Overall, OOD generalization is driven primarily by removing absolute-value channels, whereas a relative $(r,c)$ bias alone provides limited benefit when absolute shortcuts remain.



\begin{table}[t]
	\centering
\caption{\textbf{Ablation: positional encodings.}
Removing PE improves OOD bag-level accuracy while keeping relative biases fixed.}
	\label{tab:ablation_pe_hurt_ours}
	\begin{tabular}{lcc}
		\toprule
		\textbf{Variant (Ours)} &
		$\mathrm{Acc}_{\text{bag}}(D_{\text{test1}})$ &
		$\mathrm{Acc}_{\text{bag}}(D_{\text{test2}})$ \\
		\midrule
		Ours + PE & 94.41(2.81) & 89.44(4.09) \\
		\textbf{Ours} & \textbf{98.47(0.35)} & \textbf{98.22(0.48)} \\
		\bottomrule
	\end{tabular}
\end{table}


\begin{table}[t]
	\centering
\caption{\textbf{Ablation: relative spatial bias alone.}
Under a standard Transformer with PE$(r,c)$, replacing absolute attention with a relative $(r,c)$ bias provides limited OOD improvement.}


	\label{tab:ablation_rel_rc}
	\begin{tabular}{lcc}
		\toprule
		\textbf{Variant} &
		$\mathrm{Acc}_{\text{bag}}(D_{\text{test1}})$ &
		$\mathrm{Acc}_{\text{bag}}(D_{\text{test2}})$ \\
		\midrule
		Transformer & 91.67(1.13) & 86.01(3.07) \\
		Relative $(r,c)$ & 92.75(4.85) & 87.52(2.69) \\
		\bottomrule
	\end{tabular}
\end{table}

\subsection{Instance-Level Coherence on Pieri Cases}
\label{sec:pieri_coherence}

Our ultimate goal is rule discovery from the predicted selected set $\hat{\mathcal Z}$, but for general instances the latent selections $\{z_u\}$ are unobserved, making instance-level evaluation unavailable. We therefore use Pieri cases, where a proved rule provides a reference set $\mathcal Z_{\text{Pieri}}$, as a qualitative anchor. Exact set agreement is not the primary objective: even within Pieri regimes, valid instance-level selections need not be unique, and different selections can induce the same aggregate signed counts. Despite not being optimized to match Pieri targets, we observe \textbf{100\% inclusion} and \textbf{66.08\% exact-set match}, indicating broad consistency with classical Pieri structure.
This is especially informative because Pieri instances appear in the OOD split $D_{\text{test2}}$: although the selector is trained only on small regimes, it still produces Pieri-consistent selections under extrapolation, suggesting that it learns a transferable rule rather than memorizing absolute patterns.



\section{Conclusion}
\label{sec:conclusion}

We proposed a differentiable framework for Latent Signed Counting that learns instance-level selection rules from integer-valued aggregate supervision. The core component, the Probabilistic Convolutional Aggregator, computes the PMF of a signed sum via convolution, enabling end-to-end maximum-likelihood training.

We demonstrated the approach on a challenging mathematical rule-discovery task, the conjectural $(q,t)$-Littlewood--Richardson factorization problem, and demonstrated improved extrapolation with the Purely-Relative Transformer, which removes absolute positional shortcuts through relative attention biases with window clamping. Although motivated by this mathematical setting, the methodology is domain-agnostic; identifying and validating applications in other scientific domains with integer aggregate measurements is an important direction for future work.




\section*{Impact Statement}
This paper presents work whose goal is to advance the field of Machine
Learning. There are many potential societal consequences of our work, none
which we feel must be specifically highlighted here.


\bibliography{reference}  
\bibliographystyle{icml2026}

%%%%%%%%%%%%%%%%%%%%%%%%%%%%%%%%%%%%%%%%%%%%%%%%%%%%%%%%%%%%%%%%%%%%%%%%%%%%%%%
%%%%%%%%%%%%%%%%%%%%%%%%%%%%%%%%%%%%%%%%%%%%%%%%%%%%%%%%%%%%%%%%%%%%%%%%%%%%%%%
% APPENDIX
%%%%%%%%%%%%%%%%%%%%%%%%%%%%%%%%%%%%%%%%%%%%%%%%%%%%%%%%%%%%%%%%%%%%%%%%%%%%%%%
%%%%%%%%%%%%%%%%%%%%%%%%%%%%%%%%%%%%%%%%%%%%%%%%%%%%%%%%%%%%%%%%%%%%%%%%%%%%%%%

\end{document}

% This document was modified from the file originally made available by
% Pat Langley and Andrea Danyluk for ICML-2K. This version was created
% by Iain Murray in 2018, and modified by Alexandre Bouchard in
% 2019 and 2021 and by Csaba Szepesvari, Gang Niu and Sivan Sabato in 2022.
% Modified again in 2023 and 2024 by Sivan Sabato and Jonathan Scarlett.
% Previous contributors include Dan Roy, Lise Getoor and Tobias
% Scheffer, which was slightly modified from the 2010 version by
% Thorsten Joachims & Johannes Fuernkranz, slightly modified from the
% 2009 version by Kiri Wagstaff and Sam Roweis's 2008 version, which is
% slightly modified from Prasad Tadepalli's 2007 version which is a
% lightly changed version of the previous year's version by Andrew
% Moore, which was in turn edited from those of Kristian Kersting and
% Codrina Lauth. Alex Smola contributed to the algorithmic style files.
